\reading{CR-16}{Netting, Close-out and Related Aspects}{DEFER}{0}

\begin{crkeybox}[Learning objectives for this reading]
\begin{enumerate}[label=\textbf{\alph*.},leftmargin=2em]
\item Explain the purpose of an International Swaps and Derivatives Association (ISDA) master agreement.
\item Summarize netting and close-out procedures (including multilateral netting), explain their advantages and disadvantages, and describe how they fit into the framework of the ISDA master agreement.
\item Describe the effectiveness of netting in reducing credit exposure under various scenarios.
\item Describe the mechanics of termination provisions and trade compressions and explain their advantages and disadvantages.
\item Provide examples of trade compression of derivative positions, calculate net notional exposure amount, and identify the party holding the net contract position in a trade compression.
\item Identify and describe termination events and discuss their potential effects on parties to a transaction.
\end{enumerate}
\end{crkeybox}

% =====================================================================
\lo{CR-16 a}{Explain the purpose of an International Swaps and Derivatives Association (ISDA) master agreement.}

The ISDA Master Agreement standardizes over-the-counter agreements to reduce
legal uncertainty and mitigate credit risk. This is accomplished by creating a
framework that specifies OTC agreement terms and conditions related to
\term{collateral}, \term{netting}, and \term{termination events}. The Master
Agreement can cover multiple transactions by forming a single legal contract with
an indefinite term.

The features and the eight default events it covers are set out at
\textbf{CR-14 e}.

% =====================================================================
\lo{CR-16 b}{Summarize netting and close-out procedures (including multilateral netting), explain their advantages and disadvantages, and describe how they fit into the framework of the ISDA master agreement.}

\subsection*{1) Netting}

Netting combines cash flows from different contracts with a single counterparty
into a single net amount, reducing counterparty risk and enhancing operational
efficiency. It is useful for offsetting exposures, especially when positions move
in opposite directions or when unwinding trades.

\term{Payment netting} consolidates cash flows, reducing the number of individual
settlements and so reducing settlement risk.

\begin{center}
\footnotesize
\begin{tabular}{L{7.0cm} L{7.5cm}}
\toprule
\thead{Advantages of netting} & \thead{Disadvantages of netting} \\
\midrule
\term{Exposure reduction}: offsets exposures, reducing risk and improving
operational efficiency &
  Netted exposures can be \textbf{volatile and difficult to control} \\
\term{Removal of market risk}: helps unwind positions with counterparties &
  A counterparty exiting a trade may get \textbf{less favorable terms} for the
  offsetting transaction \\
\term{Counterparty risk reduction}: trading multiple positions with the same
counterparty reduces counterparty risk and can bring more favorable terms & \\
\term{Enhanced stability}: reduces the motivation to terminate contracts with
troubled counterparties, promoting financial stability & \\
\bottomrule
\end{tabular}
\end{center}

Without netting, exposures and risks remain \emph{additive}.

\subsection*{2) Close-out netting}

A specific type of netting agreement governed by contracts like the ISDA Master
Agreement. It comes into play upon counterparty \emph{default}, allowing
termination of contracts and offsetting of exposures. It significantly reduces
exposure to the insolvent counterparty, as the solvent party can offset positive
and negative positions.

\begin{itemize}
\item Protects the solvent entity by allowing termination and replacement of
contracts with a defaulted counterparty.
\item Limits uncertainty in the value of the solvent entity's position.
\item Enables freezing and hedging of exposures.
\end{itemize}

\begin{crtrapbox}[Close-out netting is not an acceleration clause]
Close-out netting differs from an \term{acceleration clause}, which makes
\emph{future payments immediately due} upon a credit event. Close-out netting
terminates and offsets; acceleration pulls the schedule forward.
\end{crtrapbox}

\subsection*{3) Multilateral netting}

Extends netting beyond two parties by involving a central entity like a
clearinghouse. Bilateral netting involves arrangements between two entities;
multilateral netting involves arrangements between multiple counterparties,
mitigating counterparty and operational risk.

\begin{center}
\footnotesize
\begin{tabular}{L{7.0cm} L{7.5cm}}
\toprule
\thead{Advantage} & \thead{Disadvantages of multilateral netting} \\
\midrule
Mitigates counterparty and operational risks &
  \textbf{Mutualizes} counterparty risk, reducing the incentive for individual
  credit monitoring \\
 & May lead to \textbf{redundant trading positions} \\
 & Requires \textbf{trading disclosure}, potentially compromising proprietary
  information confidentiality \\
\bottomrule
\end{tabular}
\end{center}

% =====================================================================
\lo{CR-16 c}{Describe the effectiveness of netting in reducing credit exposure under various scenarios.}

Netting reduces exposure by offsetting positive and negative mark-to-market
values of trades with the same counterparty. How much benefit it delivers depends
entirely on whether the instruments can go negative.

\begin{center}
\footnotesize
\begin{tabular}{L{5.6cm} L{8.9cm}}
\toprule
\thead{Instrument type} & \thead{Netting benefit} \\
\midrule
Instruments that can take a \textbf{negative} MtM value during their life, for
example long options without premiums &
  \textbf{Benefit netting the most.} They supply the negative values that offset
  the positives \\
Instruments with \textbf{only or mostly positive} MtM values, for example equity
options with premiums, swaptions, caps, floors, and FX options &
  \textbf{Less beneficial.} Still often included, for two reasons: future trades
  with negative MtM may offset them, and it reduces residual counterparty risk
  when unwinding positions \\
\bottomrule
\end{tabular}
\end{center}
\figcap{The discriminating property is not the instrument's riskiness but the
\emph{sign range} of its mark-to-market. An instrument that can only ever be an
asset to you has nothing to contribute to an offset.}

\subsection*{The scale of the reduction}

\begin{center}
\footnotesize
\begin{tabular}{l r r r}
\toprule
 & \thead{No netting} & \thead{Bilateral netting} & \thead{Central clearing} \\
\midrule
Counterparty 1 & 170 & 50 & 30 \\
Counterparty 2 & 90 & 20 & 0 \\
Counterparty 3 & 100 & 50 & 0 \\
CCP & --- & --- & 30 \\
\midrule
\textbf{Total} & \textbf{360} & \textbf{120} & \textbf{60} \\
\bottomrule
\end{tabular}
\end{center}

\begin{center}
\begin{tikzpicture}
\begin{axis}[
  ybar, bar width=22pt, width=11.4cm, height=5.4cm,
  ylabel={Total exposure},
  ylabel style={font=\sffamily\scriptsize}, tick label style={font=\sffamily\scriptsize},
  symbolic x coords={No netting,Bilateral netting,Central clearing},
  xtick=data, ymin=0, ymax=430,
  axis y line*=left, axis x line*=bottom,
  nodes near coords, every node near coord/.append style={font=\sffamily\scriptsize\bfseries}]
\addplot[draw=FRMNavy, fill=FRMNavy!45, area legend]
  coordinates {(No netting,360) (Bilateral netting,120) (Central clearing,60)};
\end{axis}
\end{tikzpicture}
\end{center}
\figcap{Total exposure across a three-party network under the three regimes.
Bilateral netting removes two-thirds of the exposure, and moving to central
clearing halves what is left, for an overall reduction of 83\%. The second step
is worth less than the first in absolute terms, which is easy to get backwards
when the headline is that central clearing is the stronger mechanism.}

% =====================================================================
\lo{CR-16 d}{Describe the mechanics of termination provisions and trade compressions and explain their advantages and disadvantages.}

\subsection*{Termination features}

Termination features allow early termination of trades \emph{before} counterparty
issues arise, that is, before a counterparty goes bankrupt.

\begin{enumerate}[label=\alph*)]
\item \term{Reset agreements.} Adjust trade parameters to bring a trade back
``at-the-money'' when it is heavily skewed in one direction, that is, heavily in
the money.
\item \term{Additional termination events (ATEs) or break clauses.} Allow
termination upon a decline in counterparty creditworthiness.
  \begin{itemize}
  \item \emph{Bilateral}: either party can terminate.
  \item Useful for exiting long-dated trades with declining counterparty quality.
  \item Triggered by mandatory, optional, or trigger-based events such as a
  ratings downgrade.
  \end{itemize}
\item \term{Walkaway clauses} (discouraged). Allow a party to benefit from
counterparty default by avoiding net liabilities while still claiming positive
MtM.
\end{enumerate}

\begin{crtrapbox}[Why break clauses are less used than collateral]
Break clauses are less popular than collateralization, and the reason is
behavioural rather than legal: the \term{banker's paradox}. Entities avoid
exercising the clause early in order to maintain good client relationships,
which is precisely when exercising it would be useful.

\smallskip
Walkaway clauses are uncommon for three reasons: they increase costs for the
counterparty in default, they create moral hazard, and they hide transaction
risks.
\end{crtrapbox}

\subsection*{Trade compression}

Trade compression eliminates overlapping positions across multiple counterparties
in a portfolio. It lowers the gross notional amount and the quantity of trades,
optimizing net exposure \emph{without altering the risk profile}.

\begin{enumerate}
\item Participants submit trades for compression based on their acceptable risk
levels.
\item Trades are paired with corresponding counterparties and netted to form a
single, consolidated contract.
\item The net contract coupon is determined using the \term{weighted average} of
the coupons of the original contracts.
\end{enumerate}

% =====================================================================
\lo{CR-16 e}{Provide examples of trade compression of derivative positions, calculate net notional exposure amount, and identify the party holding the net contract position in a trade compression.}

\begin{crexbox}[Trade compression on an index]
A participant holds the following positions, all on the XYZ index and all
maturing 20 March 2025 with a coupon of 100.

\smallskip
\centering\footnotesize
\begin{tabular}{l r l l r l}
\toprule
\thead{Reference entity} & \thead{Notional} & \thead{Position} &
\thead{Maturity} & \thead{Coupon} & \thead{Counterparty} \\
\midrule
XYZ index & 200 & Long & 20 March 2025 & 100 & Counterparty X \\
XYZ index & 100 & Short & 20 March 2025 & 100 & Counterparty Y \\
XYZ index & 75 & Short & 20 March 2025 & 100 & Counterparty Z \\
\bottomrule
\end{tabular}

\raggedright
\smallskip
Calculate the net notional exposure amount and identify the party holding the net
contract position.
\tcblower
\textbf{Net notional.} Long positions less short positions:
\[ 200 - 100 - 75 = 25 \text{ Long} \]

\textbf{Net contract coupon.} The weighted average of the original coupons. All
three are 100, so the net coupon is 100.

\textbf{Party holding the net contract.} \term{Counterparty X}, the counterparty
to the surviving long position.

\smallskip
\centering\footnotesize
\begin{tabular}{l r l l r l}
\toprule
\thead{Reference entity} & \thead{Notional} & \thead{Position} &
\thead{Maturity} & \thead{Coupon} & \thead{Counterparty} \\
\midrule
\textbf{Net contract} & \textbf{25} & \textbf{Long} & 20 March 2025 & 100 &
\textbf{Counterparty X} \\
\bottomrule
\end{tabular}

\raggedright
\smallskip
\textbf{What compression achieved.} Gross notional falls from
$200+100+75 = 375$ to 25, a reduction of 350 or \textbf{93.3\%}, and three trades
become one. The net risk position is unchanged throughout.
\end{crexbox}

\begin{crtrapbox}[Gross versus net]
Compression reduces the \emph{gross} notional and the trade count. It does not
change the net position, which was 25 long before compression and 25 long after.
The point of the exercise is operational and capital efficiency, not risk
reduction.
\end{crtrapbox}

% =====================================================================
\lo{CR-16 f}{Identify and describe termination events and discuss their potential effects on parties to a transaction.}

The termination events themselves are the ATEs and break clauses set out at
objective d, triggered by mandatory, optional, or trigger-based events such as a
ratings downgrade. Their effects on the parties are:

\begin{itemize}
\item \term{On the terminating party}: an exit from a long-dated trade with a
deteriorating counterparty, before default occurs.
\item \term{On the counterparty}: a forced unwind at a potentially unfavourable
moment, since the trigger is by construction a point of weakness.
\item \term{On both}: under a bilateral clause, either side can invoke it, so the
protection is symmetric even where the credit deterioration is not.
\end{itemize}

The eight default events that the ISDA Master Agreement itself covers are listed
at \textbf{CR-14 e}.
